\chapter{Vyhodnotenie}
\label{kap:vyhodnotenie}

V tejto kapitole sa zameriame na vyhodnotenie výsledkov, ktoré sme dosiahli pri vývoji a implementácii autonómneho agenta pre hru Trackmania. Na základe teoretických východísk prezentovaných v kapitole \ref{kap:teoria} a praktickej implementácie v kapitole \ref{kap:prakticka_cast}, sme otestovali nášho agenta na troch rôznych tratiach, na ktorých doposial nebol trénovaný. Výkony nášho agenta sme porovnali s dvoma konkurenčnými projektami: Trackmania AI a TMRL. V nasledujúcich sekciách detailne rozoberieme výsledky a analyzujeme rýchlosť odozvy jednotlivých systémov.

Testovanie prebehlo na troch rôznych tratiach, ktoré vytvoril autor projektu Trackmania AI aj s časmi pre jeho agenta. Pre nášho agenta tieto trate poskytujú dostatočnú rozmanitosť na overenie schopnosti agenta adaptovať sa na nové a neznáme prostredie.

\section{Porovnanie časov prejdenia trate}
\label{sec:cas_prejdenia}

Porovnali sme priemerný a najlepší čas prejdenia trate naším agentom, Trackmania AI agentom a TMRL agentom. Výsledky sú uvedené v tabuľkách \ref{tab:cas} a \ref{tab:najlepsi_cas}. Môžeme si všimnúť, že agent TMRL má výrazný problém prejsť mapu v nočnom čase, keďže sa spoliehal na spracovanie obrazu.

\begin{table}[h]
\centering
\begin{tabular}{|c|c|c|c|}
\hline
\textbf{Trať} & \textbf{Náš agent (s)} & \textbf{Trackmania AI (s)} & \textbf{TMRL (s)} \\ \hline
\textbf{Test2} & 16.258 & 14.655 & 15.367 \\ \hline
\textbf{Nascar2} & 19.227 & 17.791 & 20.090 \\ \hline
\textbf{Night} & 11.937 & 10.999 & 27.458 \\ \hline
\end{tabular}
\caption{Porovnanie priemerného času prejdenia trate}
\label{tab:cas}
\end{table}

\begin{table}[h]
\centering
\begin{tabular}{|c|c|c|c|}
\hline
\textbf{Trať} & \textbf{Náš agent (s)} & \textbf{Trackmania AI (s)} & \textbf{TMRL (s)} \\ \hline
\textbf{Test2} & 15.150 & 14.109 & 14.824 \\ \hline
\textbf{Nascar2} & 18.143 & 16.180 & 17.650 \\ \hline
\textbf{Night} & 11.739 & 10.576 & 18.940 \\ \hline
\end{tabular}
\caption{Porovnanie najlepšieho času prejdenia trate}
\label{tab:najlepsi_cas}
\end{table}

\section{Porovnanie rýchlosti odozvy}
\label{sec:odozva}

Rýchlosť odozvy je kľúčovým faktorom pri riadení autonómnych vozidiel, najmä v dynamických prostrediach ako je Trackmania. Nasledujúca tabuľka zobrazuje porovnanie rýchlosti odozvy medzi našim agentom, Trackmania AI agentom a TMRL agentom.

\begin{table}[h]
\centering
\begin{tabular}{|c|c|c|c|}
\hline
\textbf{Veľkosť trate} & \textbf{Náš agent (ms)} & \textbf{Trackmania AI (ms)} & \textbf{TMRL (ms)} \\ \hline
\textbf{Malá trať} & 13 & 13 & 75 \\ \hline
\textbf{Stredná trať} & 20 & 15 & 75 \\ \hline
\textbf{Veľká trať} & 34 & 23 & 75 \\ \hline
\end{tabular}
\caption{Porovnanie rýchlosti odozvy agenta}
\label{tab:odozva}
\end{table}

\section{Záverečné zhodnotenie}
\label{sec:zaver}

Výsledky testovania ukázali, že náš agent dosiahol konkurenčný výkon v porovnaní s projektom Trackmania AI a TMRL. Priemerný a najlepší čas prejdenia tratí bol porovnateľný s týmito projektmi, pričom náš agent vykazoval vyššiu konzistenciu v rôznych podmienkach a rýchlejšiu odozvu na dynamické zmeny v prostredí.

Náš agent použil iný prístup ako ostatné projekty, pričom dosiahol podobné výsledky pri prejdení tratí. Hlavné prednosti nášho agenta sú:

\begin{itemize}
    \item \textbf{Plná autonómnosť:} Náš agent nepotrebuje žiadnu ľudskú interakciu na prejdenie trate.
    \item \textbf{Rýchla odozva:} Náš agent má výrazne rýchlejšiu odozvu v porovnaní s konkurenčnými projektmi.
    \item \textbf{Flexibilita do budúcnosti:} Náš agent nie je limitovaný na 2D priestor, čo nám umožňuje rozšíriť jeho schopnosti aj na trate s výškovým prevýšením alebo slučkami.
\end{itemize}

Tieto výsledky potvrdzujú efektívnosť nášho prístupu a potenciál pre ďalšie zlepšovanie v smere autonómneho riadenia v hre Trackmania.
